\section{Review Boundaries and Journal Projection Map}
\label{sec:oc133-review-boundaries-journal-map}

This chapter translates likely reviewer objections into the form used by the
release. Each objection threatens a particular claim, calls for a particular
kind of evidence, leaves a residual risk, and has a reopening condition. The
purpose is not to win by phrasing; it is to make criticism precise enough that
the manuscript can be repaired when the criticism is correct.

\subsection{The tuple may be only vocabulary}

The threatened claim is that \(K=(\Omega,\partial\Omega,A,\Theta,P,J,C,k,M)\)
is a useful typed model object rather than a decorative list. The answer is the
component discipline: each component names a distinct review question about
state, boundary, observation, threshold, potential, flow, cycle, persistence,
or context. The residual risk is redundancy. The reopening condition is a case
where a component can be removed without changing any verdict, witness,
boundary, or falsifier.

\subsection{Boundaries may be metaphorical}

The threatened claim is that \(\partial\Omega\) can do scientific work. The
answer is to read boundaries as typed admissibility conditions: they separate
states, events, or relations that preserve the declared continuum from those
that demote, split, kill, or require a narrower claim. The residual risk is that
some domains need richer boundary geometry. The reopening condition is a
declared boundary that fails to separate admissible and inadmissible states in
the cited case.

\subsection{Continuumness may be underdefined}

The threatened claim is that \(k\) names persistence through change. The answer
is that \(k\) is always read together with state region, boundary, cycles, and
context; it is not a mystical liveness score. The residual risk is that
empirical proxies for liveness remain domain-bound. The reopening condition is
a live/dead distinction that cannot be reproduced under the stated assumptions.

\subsection{K-levels may be labels}

The threatened claim is that K-levels are nested witness strata. The answer is
the retained-witness and demotion discipline: an alleged level must add a
constraint, observable consequence, or failure criterion. The residual risk is
that some didactic examples are ahead of domain validation. The reopening
condition is an adjacent K transition that adds no retained witness.

\subsection{Operator claims may be too broad}

The threatened claim is that operators preserve type discipline across updates.
The answer is to keep operator claims within declared update families and
assumption classes. The residual risk is domain adapter complexity. The
reopening condition is an operator that silently changes the claim type without
a declared rule.

\subsection{Formal support may not cover prose}

The threatened claim is that proof, finite witnesses, and formalization
inventory support the public wording. The answer is alignment: the public
sentence must stay within the exact support class. The residual risk is partial
mechanization. The reopening condition is a dependency, assumption, or
formalization mismatch.

\subsection{Replay rows may be overread}

The threatened claim is that replay examples provide bounded support. The
answer is the replay evidence pattern: source snapshot, formula, residual,
comparator, negative control, and falsifier. The residual risk is that a row is
illustrative rather than validation-grade. The reopening condition is any
missing or failed replay component.

\subsection{Prior art may absorb the delta}

The threatened claim is that OC contributes residual structure beyond existing
systems theory, cybernetics, autopoiesis, dynamical systems, complexity science,
formal verification, formal ontology, and compositional modeling. The answer is
the comparator matrix. The residual risk is real: a stronger comparator can
narrow OC. The reopening condition is a comparator that explains the same row
with less burden.

\subsection{Figures and tables may imply too much}

The threatened claim is that visuals and tables are argument surfaces. The
answer is that each reader-facing figure or table must have a semantic role,
caption, formula or evidence anchor, and rendered-layout check. The residual
risk is visual compression. The reopening condition is a figure or table that
implies a claim without its support anchor.

\subsection{Journal projections may drift}

The threatened claim is that venue-specific article packages remain projections
from the same release. The answer is a no-send journal package discipline:
venue requirements, component manifests, source maps, and disclosure statements
are generated from the release package. The residual risk is loss of nuance
during shortening. The reopening condition is a projection that introduces a
claim not present in the release or drops a necessary limitation.

\subsection{Foundations-of-science review stance}

For a foundations-oriented venue, the central editorial question is whether the
manuscript has a real conceptual object and a serious relation to prior systems
thought. The projection should foreground the tuple, the continuity problem,
the prior-art comparator matrix, and the limits of present proof. It should not
pretend that the current release has completed every empirical domain. A
foundations reviewer can therefore attack the model grammar, the identity
conditions, the ontology of boundaries, and the relation to existing systems
theory.

\subsection{Physics-oriented review stance}

For a physics-oriented venue, the central editorial question is whether the
formal language respects the difference between mathematical structure and
physical validation. The projection should foreground state space, boundary,
threshold, flow, finite witness, and falsifier discipline, while treating
high-level numeric ranges as illustrative unless they have row-level source and
derivation. A physics reviewer can therefore attack dimensional consistency,
observable definition, benchmark selection, and whether any physical claim has
been promoted beyond its evidence.

\subsection{Biotheory and systems-biology stance}

For a biotheory venue, the central question is whether continuumness, liveness,
residue, rebirth, and identity are operational rather than metaphorical. The
projection should foreground lifecycle distinctions, cycles, maintenance,
boundary crossing, and demotion conditions. A reviewer can ask whether the
model distinguishes organism, trace, residue, population renewal, and identity
continuation without smuggling metaphysics into biological evidence.

\subsection{Management and enterprise stance}

For a management, flexible-systems, or enterprise-architecture venue, the
central question is whether the model helps decision-makers inspect complex
systems without losing scientific discipline. The projection should foreground
domain projection, enterprise examples, AI and security routes, reader-facing
tables, and limitations. A reviewer can ask whether the model improves
architecture reasoning, whether it identifies falsifiers, and whether it avoids
turning strategy language into unsupported universals.

\subsection{Computational-biology and reproducibility stance}

For a computational or reproducibility venue, the central question is whether
the package is inspectable. The projection should foreground source snapshots,
commands, retrospective replay rows, finite checks, table layout, data/code statements,
and failure interpretation. A reviewer can ask whether an independent reader
can reproduce the bounded check, whether hashes bind to the right files, and
whether the claim wording changes when a replay row fails.

\subsection{Chemistry and interdisciplinary-data stance}

For an interdisciplinary data venue, the central question is whether examples,
tables, and source files are sufficiently explicit for external reuse. The
projection should foreground the data/code manifest, evidence anchors, SI
boundary, visual/table quality, and conservative interpretation of replay
examples. A reviewer can ask whether each table has a source, whether each
figure demonstrates a claim, and whether domain examples are demoted when
source or comparator support is incomplete.
